\section{Evaluation}\label{sec:evaluation}

We ask three questions. \textbf{RQ1}: which mistakes does each rung catch, and which does it
miss? \textbf{RQ2}: what do the added rungs cost in size and reading? \textbf{RQ3}: do the
declared and enforced rungs change what coding agents do? Every number below can be regenerated
from the repository's \repo{evaluation/} directory \citep{taskboardrepo}.

\subsection{RQ1: violation seeding}

A seeding script plants one canonical mistake at a time in a receiver, runs its single verify
command, records the earliest mechanism that objects, and restores the file. It refuses to start
unless the receiver is green, and checks that it is green again at the end. The mechanisms, in
order of how early they give feedback, are the compiler, an architecture rule, the requirement
check, a behavioural test, and nothing. Table~\ref{tab:violations} lists the fourteen mistakes:
twelve structural, and two that break the trace between requirements and scenarios.

\begin{table*}[t]
  \centering\small
  \caption{The fourteen seeded mistakes, and the mechanism that caught each in Java / Python /
  TypeScript now. C compiler, A architecture rule, R requirement check, B behavioural test.}
  \label{tab:violations}
  \begin{tabularx}{\linewidth}{l|X|c}
    \hline
    & \textbf{Mistake} & \textbf{J / P / T} \\
    \hline\hline
    V1 & The API bypasses the service and uses a repository & A / A / A \\
    \hline
    V2 & A service signals failure with an HTTP concept & A / A / B \\
    \hline
    V3 & A service issues SQL & A / A / A \\
    \hline
    V4 & A component constructs its own collaborator & A / A / A \\
    \hline
    V5 & A service reads the environment & A / A / A \\
    \hline
    V6 & A new domain error has no HTTP status & C / A / C \\
    \hline
    V7 & A route catches a domain error and answers ad hoc & B / A / A \\
    \hline
    V8 & A side effect is inlined in the service & A / A / A \\
    \hline
    V9 & The test double drifts from the real repository & B / B / B \\
    \hline
    V10 & A business rule is placed in the route & A / A / A \\
    \hline
    V11 & A setting is hard-coded in the service & B / B / B \\
    \hline
    V12 & A framework type enters the domain & A / A / A \\
    \hline
    V13 & A requirement has no scenario & R / R / R \\
    \hline
    V14 & A scenario cites no requirement & R / R / R \\
    \hline
  \end{tabularx}
\end{table*}

Figure~\ref{fig:seeding} compares the receivers at commit \texttt{8c4a74a}, before the work
described here, with the receivers now. Before, the single gate caught 30 of 36 structural
mistakes. The six misses were the same two mistakes in every language: V8, a side effect inlined
in the service, and V10, a business rule placed in the route. That the gap was identical across
three languages and three enforcement styles is itself a result: it belonged to the capability,
not to any receiver. Both mistakes preserve behaviour, so no behavioural test could have caught
them; only a structural rule could. After two such rules were added, and the requirement check
with its two seeds, the gate catches all 42 mistakes. The new routes rule also moved V7 from a
behavioural test to an architecture rule in Python and TypeScript, which is earlier feedback for
the same mistake. Application code did not change between the two runs; only tests, manifests
and requirement files did.

\begin{figure*}[t]
  \centering
  % Seeding results, before (8c4a74a, 12 structural seeds) and now (14 seeds), per receiver.
\begin{tikzpicture}
  \begin{axis}[
      xbar stacked,
      width=0.86\linewidth,
      height=6.2cm,
      bar width=9pt,
      xmin=0, xmax=14,
      xtick={0,2,...,14},
      xlabel={seeded mistakes},
      ytick=data,
      yticklabels={{TypeScript, now},{TypeScript, then},{Python, now},{Python, then},{Java, now},{Java, then}},
      yticklabel style={font=\footnotesize},
      tick label style={font=\footnotesize},
      label style={font=\footnotesize},
      legend style={at={(0.5,-0.24)}, anchor=north, legend columns=5, font=\footnotesize, draw=none},
      enlarge y limits=0.12,
      axis x line*=bottom,
      axis y line*=left,
      xmajorgrids, grid style={muted!20},
    ]
    % order of rows (bottom to top): Java then, Java now, Python then, Python now, TS then, TS now
    \addplot[fill=compiler, draw=none] coordinates {(1,5) (1,4) (0,3) (0,2) (1,1) (1,0)};
    \addplot[fill=caught, draw=none]   coordinates {(6,5) (8,4) (7,3) (10,2) (5,1) (8,0)};
    \addplot[fill=reqc, draw=none]     coordinates {(0,5) (2,4) (0,3) (2,2) (0,1) (2,0)};
    \addplot[fill=behav, draw=none]    coordinates {(3,5) (3,4) (3,3) (2,2) (4,1) (3,0)};
    \addplot[fill=missed, draw=none]   coordinates {(2,5) (0,4) (2,3) (0,2) (2,1) (0,0)};
    \legend{compiler, architecture rule, requirement check, behavioural test, not caught}
  \end{axis}
\end{tikzpicture}

  \caption{Seeded mistakes by the earliest mechanism that caught them, per receiver, at commit
  \texttt{8c4a74a} (12 structural seeds) and now (14 seeds).}
  \label{fig:seeding}
\end{figure*}

\subsection{RQ2: cost}

Table~\ref{tab:size} reports size. Application code is unchanged. Test code grew by between 31
and 89 lines per receiver, and each manifest by between 80 and 83 words, remaining under the 500
words the plan allows. The added rungs therefore cost about one screen of test code and one
paragraph of prose per receiver, most of it the requirement check and the ids on existing tests.

\begin{table*}[t]
  \centering\small
  \caption{Size of the receivers at commit \texttt{8c4a74a} and now (\repo{evaluation/metrics*.json}).}
  \label{tab:size}
  \begin{tabular}{l|r|r|r|r|r|r}
    \hline
    \multicolumn{1}{l|}{} & \multicolumn{2}{c|}{\textbf{Application lines}} & \multicolumn{2}{c|}{\textbf{Test lines}} & \multicolumn{2}{c}{\textbf{Manifest words}} \\
    \cline{2-3}\cline{4-5}\cline{6-7}
    \textbf{Receiver} & \textbf{then} & \textbf{now} & \textbf{then} & \textbf{now} & \textbf{then} & \textbf{now} \\
    \hline\hline
    Java       & 283 & 283 & 287 & 376 & 397 & 480 \\
    \hline
    Python     & 223 & 223 & 193 & 254 & 347 & 427 \\
    \hline
    TypeScript & 242 & 242 & 193 & 224 & 388 & 468 \\
    \hline
  \end{tabular}
\end{table*}

\subsection{RQ3: a pilot with coding agents}

Eighteen coding agents, each starting with no memory, implemented the same feature request, task
deletion, in copies of the receivers at commit \texttt{8c4a74a}. Each copy offered one, two or
three of the carriers: the exemplar only (L1), the exemplar and the manifest (L1--L2), or the
repository as published (L1--L3). The design crossed three languages, three conditions and two
models, with one run per cell. A hidden acceptance test of four cases, never shown to an agent,
judged the result, and a scripted audit fixed in advance of inspection recorded where each agent
placed each part of the change.

\begin{table*}[t]
  \centering\small
  \caption{Pilot results by condition (six runs each; \repo{evaluation/agent\_experiment/}).}
  \label{tab:pilot}
  \begin{tabular}{l|c|c|c|c}
    \hline
    \textbf{Condition} & \textbf{Acceptance passed} & \textbf{Departures from convention} & \textbf{Mean tokens} & \textbf{Mean seconds} \\
    \hline\hline
    Exemplar (L1)            & 6 / 6 & 2 & 61\,650 & 121 \\
    \hline
    + Manifest (L1--L2)      & 6 / 6 & 0 & 61\,990 & 132 \\
    \hline
    + Enforced (L1--L3)      & 6 / 6 & 0 & 65\,696 & 151 \\
    \hline
  \end{tabular}
\end{table*}

All eighteen runs delivered a passing suite and passed every acceptance case
(Table~\ref{tab:pilot}). Two departed from the conventions, both by the smaller model in the
exemplar-only condition. One, in Java, logged the deletion inside the service instead of
publishing an event: the V8 mistake, which no rule caught at the time. The deletion request's
last clause, ``whenever a task is deleted the team must be told'', is an event-driven requirement
in EARS terms, and the mapping in Table~\ref{tab:ears} places it in an event handler. The other,
in Python, did not add the new operation to the in-memory double, which only the type checker in
the full rule set notices. With one run per cell these results are suggestive, not significant.
The pilot has not yet been repeated against the receivers as they stand now, with the deletion
request written in EARS; that is the next experiment.
